\documentclass[../main.tex]{subfiles}
\begin{document}

% capitulos/05_resultados.tex

\chapter{Resultados}
\label{C:resultados}

\section{Descripción del experimento ejecutado}
\label{sec:experimento-ejecutado}

El experimento completo se ejecutó el 24~de julio de 2026 sobre el \emph{host}
\texttt{pc-pantu}. El \emph{runner} completó las 240~ejecuciones en aproximadamente
cinco horas y media sin ningún error ni ejecución parcial. El fichero CSV resultante
contiene exactamente 240~filas, veinte por cada una de las doce celdas del factorial
($4~\text{escenarios} \times 3~\text{estrategias} \times 20~\text{repeticiones}$),
ejecutadas en orden aleatorio con semilla fija para evitar que la estrategia quede
confundida con el momento de ejecución.

A diferencia de versiones preliminares del laboratorio, el tiempo de detección
($\tdetect$) se mide ahora a partir del \emph{timestamp} real que el \emph{indexer}
asigna a la primera alerta del grupo del escenario, y no del instante en que el
sondeo la descubre. El sondeo periódico cada 5~s solo introduce latencia en la
ejecución, no en la medida: el valor registrado es el que el \gls{siem} asignó al
generar la alerta. Antes de cada ataque se espera además a que el ruido de eventos
\gls{fim} generado por el \emph{reset} se drene por completo, con el mismo criterio
para las tres estrategias.

La verificación de la integridad del CSV, realizada mediante el \emph{script}
\texttt{01\_analisis\_estadistico.py}, confirmó que no hay valores faltantes en
ninguna columna obligatoria ni filas con estados inconsistentes.

\section{Tasa de recuperación}

La recuperación solo tiene sentido en los escenarios que provocan pérdida de datos
(E1, cifrado local; y E2, destrucción de copias). Los escenarios E3 (movimiento
lateral) y E4 (exfiltración) no cifran ni destruyen los datos del \emph{fileserver},
de modo que su tasa de recuperación no es interpretable y se excluye de este
análisis; en ellos la variable de interés es la detección. La tabla~\ref{tab:tasa-recuperacion}
presenta la tasa de recuperación exitosa (\texttt{recovery\_status = success}) por
estrategia sobre las 40~ejecuciones con pérdida de datos de cada una.

\begin{table}[ht]
  \centering
  \caption{Tasa de recuperación por estrategia (40~ejecuciones con pérdida de datos,
           escenarios E1 y E2, por estrategia).}
  \label{tab:tasa-recuperacion}
  \rowcolors{2}{tabla-par}{white}
  \begin{tabularx}{\textwidth}{X C{2.2cm} C{2cm} C{2.4cm} C{2.4cm}}
    \toprule
    \rowcolor{tabla-cabecera}
    \textcolor{white}{\textbf{Estrategia}} &
    \textcolor{white}{\textbf{Ejec.}} &
    \textcolor{white}{\textbf{Success}} &
    \textcolor{white}{\textbf{Impossible}} &
    \textcolor{white}{\textbf{Tasa de éxito}} \\
    \midrule
    A — Sin backup              & 40 &  0 & 40 &   0\,\% \\
    B — Restic + SFTP           & 40 & 20 & 20 &  50\,\% \\
    C — MinIO Object Lock       & 40 & 40 &  0 & 100\,\% \\
    \bottomrule
  \end{tabularx}
\end{table}

La diferencia entre B y C se explica íntegramente por las veinte ejecuciones del
escenario~E2: cuando el atacante destruye el repositorio \gls{sftp}, la
estrategia~B resulta irrecuperable en el 100\,\% de los casos, mientras que la
estrategia~C mantiene disponibilidad total gracias a la combinación de Object Lock
COMPLIANCE y la política de mínimo privilegio \gls{iam}. En el escenario~E1, donde
el repositorio de \emph{backup} no es objetivo del ataque, ambas estrategias alcanzan
el 100\,\% de éxito.

\begin{figure}[ht]
  \centering
  \includegraphics[width=0.85\textwidth]{figura2_tasa_recuperacion}
  \caption{Tasa de recuperación exitosa por estrategia de \textit{backup}, sobre los
           escenarios con pérdida de datos (E1 y E2). La estrategia~C alcanza el
           100\,\% incluso ante la destrucción activa del repositorio.}
  \label{fig:tasa-rec}
\end{figure}

\section{Tiempos de detección (\texorpdfstring{$\tdetect$}{T\_detect})}

La tabla~\ref{tab:tdetect} muestra la mediana del tiempo de detección por celda
experimental, junto con el rango intercuartílico [P25--P75].

\begin{table}[ht]
  \centering
  \caption{Estadísticos de $\tdetect$ por celda experimental (segundos).}
  \label{tab:tdetect}
  \rowcolors{2}{tabla-par}{white}
  \begin{tabularx}{\textwidth}{C{1.8cm} C{1cm} S[table-format=2.1] S[table-format=2.1] S[table-format=2.1]}
    \toprule
    \rowcolor{tabla-cabecera}
    \textcolor{white}{\textbf{Celda}} &
    \textcolor{white}{\textbf{$n$}} &
    \textcolor{white}{\textbf{Mediana (s)}} &
    \textcolor{white}{\textbf{P25 (s)}} &
    \textcolor{white}{\textbf{P75 (s)}} \\
    \midrule
    E1$\times$A & 20 &  0.4 &  0.4 &  0.4 \\
    E1$\times$B & 20 &  0.4 &  0.4 &  0.4 \\
    E1$\times$C & 20 &  0.5 &  0.4 &  0.5 \\
    E2$\times$A & 20 & 11.2 & 11.0 & 11.6 \\
    E2$\times$B & 20 &  0.7 &  0.7 &  0.8 \\
    E2$\times$C & 20 & 11.1 & 10.3 & 11.6 \\
    E3$\times$A & 20 &  1.5 &  0.9 &  1.8 \\
    E3$\times$B & 20 &  1.6 &  0.8 &  2.1 \\
    E3$\times$C & 20 &  1.4 &  1.2 &  2.0 \\
    E4$\times$A & 20 &  2.0 &  1.6 &  2.9 \\
    E4$\times$B & 20 &  1.2 &  1.2 &  1.2 \\
    E4$\times$C & 20 &  1.2 &  1.2 &  1.2 \\
    \bottomrule
  \end{tabularx}
\end{table}

Los tiempos de detección son consistentemente bajos, con medianas que oscilan entre
0.4\,s (E1) y 11.2\,s (E2$\times$A). La variabilidad dentro de cada celda es muy
reducida. Todas las celdas quedan muy por debajo del umbral de 60~s establecido en
H2.

\begin{figure}[ht]
  \centering
  \includegraphics[width=\textwidth]{figura3_tdetect}
  \caption{Tiempo de detección ($\tdetect$) mediano por celda experimental.
           Barras de error = IQR [P25--P75]. La línea roja discontinua
           indica el umbral de 60\,s de H2; todas las celdas quedan muy por
           debajo.}
  \label{fig:tdetect}
\end{figure}

\label{sec:tdetect-e4}
El escenario más lento es E2, y la razón es ilustrativa del mecanismo de detección.
En las estrategias A y C la alerta de E2 llega en torno a 11~s: no existe repositorio
\gls{sftp} que borrar (la estrategia~A no realiza copia y la~C la deposita en
\gls{s3}/MinIO), por lo que la detección procede del evento estructurado que confirma
el bloqueo de Object Lock, emitido tras el intento de borrado del \emph{bucket}
protegido. En la estrategia~B, en cambio, sí hay un repositorio SFTP que el atacante
borra masivamente, lo que dispara la regla de \gls{fim} sobre \texttt{/srv/backups}
en menos de un segundo (mediana 0.7~s). El escenario E4 (exfiltración), que en
iteraciones previas del laboratorio presentaba los tiempos más altos por un artefacto
de medición, se detecta ahora en torno a 1.2~s gracias a las reglas de sesión SSH y
de creación del archivo comprimido en \texttt{/tmp}.

\section{Tiempos de recuperación (\texorpdfstring{\acrshort{mttr}}{MTTR})}

El \gls{mttr}, definido como el tiempo desde el lanzamiento del ataque hasta la
verificación de la restauración, solo es calculable para las ejecuciones con
\texttt{recovery\_status = success}. Es una medida del \emph{tiempo total de
interrupción}, pero conviene interpretarla con cautela: incorpora la duración del
ataque y la latencia del sondeo de detección (cuantizada en pasos de 5~s), que
dominan el valor y lo hacen ruidoso. La comparación limpia de la velocidad de los
\emph{backends} se presenta por separado en la sección~\ref{sec:restauracion-pura}.
Los resultados por celda se recogen en la tabla~\ref{tab:mttr}.

\begin{table}[ht]
  \centering
  \caption{Estadísticos de \gls{mttr} por celda experimental (segundos).}
  \label{tab:mttr}
  \rowcolors{2}{tabla-par}{white}
  \begin{tabularx}{\textwidth}{C{1.8cm} C{0.8cm} S[table-format=3.1] S[table-format=2.1] S[table-format=3.1] l}
    \toprule
    \rowcolor{tabla-cabecera}
    \textcolor{white}{\textbf{Celda}} &
    \textcolor{white}{\textbf{$n$}} &
    \textcolor{white}{\textbf{Med. (s)}} &
    \textcolor{white}{\textbf{P25}} &
    \textcolor{white}{\textbf{P75}} &
    \textcolor{white}{\textbf{Status}} \\
    \midrule
    E1$\times$A &  0 & {---} & {---} & {---}  & impossible \\
    E1$\times$B & 20 &  12.8 &  12.8 &  12.9  & success \\
    E1$\times$C & 20 &  12.1 &   7.0 &  12.1  & success \\
    E2$\times$A &  0 & {---} & {---} & {---}  & impossible \\
    E2$\times$B &  0 & {---} & {---} & {---}  & impossible \\
    E2$\times$C & 20 &  16.8 &  16.4 &  20.9  & success \\
    E3$\times$A &  0 & {---} & {---} & {---}  & impossible \\
    E3$\times$B & 20 &   7.9 &   7.9 &  12.9  & success \\
    E3$\times$C & 20 &  12.1 &  12.0 &  12.1  & success \\
    E4$\times$A &  0 & {---} & {---} & {---}  & impossible \\
    E4$\times$B & 20 &  10.7 &  10.6 &  10.7  & success \\
    E4$\times$C & 20 &  14.8 &   9.8 &  14.9  & success \\
    \bottomrule
  \end{tabularx}
\end{table}

\begin{figure}[ht]
  \centering
  \includegraphics[width=\textwidth]{figura1_mttr_boxplot}
  \caption{Distribución del \gls{mttr} por escenario y estrategia de backup.
           El escenario~E2 no muestra datos para B porque todos los casos
           resultaron en recuperación imposible.}
  \label{fig:mttr-boxplot}
\end{figure}

\begin{figure}[ht]
  \centering
  \includegraphics[width=\textwidth]{figura4_mttr_scatter}
  \caption{Variabilidad del \gls{mttr} por repetición. Cada punto representa
           una ejecución experimental. Las líneas discontinuas indican la mediana.}
  \label{fig:mttr-scatter}
\end{figure}

\section{Velocidad de restauración}
\label{sec:restauracion-pura}

Para comparar la velocidad de los dos \emph{backends} de forma aislada se define el
\emph{tiempo de restauración puro} como el intervalo entre el inicio de la acción de
recuperación y la finalización de la restauración ($T_{restore} - T_{first\_action}$),
excluyendo el ataque y la detección. Esta métrica se mide sobre los escenarios en los
que ambas estrategias llegan a restaurar (E1, E3 y E4). La tabla~\ref{tab:restauracion}
muestra sus estadísticos por celda.

\begin{table}[ht]
  \centering
  \caption{Tiempo de restauración puro ($T_{restore}-T_{first\_action}$) por celda
           (segundos).}
  \label{tab:restauracion}
  \rowcolors{2}{tabla-par}{white}
  \begin{tabularx}{\textwidth}{C{1.8cm} C{1cm} S[table-format=1.2] S[table-format=1.2] S[table-format=1.2] X}
    \toprule
    \rowcolor{tabla-cabecera}
    \textcolor{white}{\textbf{Celda}} &
    \textcolor{white}{\textbf{$n$}} &
    \textcolor{white}{\textbf{Med. (s)}} &
    \textcolor{white}{\textbf{P25}} &
    \textcolor{white}{\textbf{P75}} &
    \textcolor{white}{\textbf{Backend}} \\
    \midrule
    E1$\times$B & 20 & 2.20 & 2.13 & 2.22 & Restic + SFTP \\
    E1$\times$C & 20 & 1.35 & 1.35 & 1.36 & MinIO (S3) \\
    E3$\times$B & 20 & 2.22 & 2.18 & 2.23 & Restic + SFTP \\
    E3$\times$C & 20 & 1.34 & 1.33 & 1.34 & MinIO (S3) \\
    E4$\times$B & 20 & 2.22 & 2.19 & 2.23 & Restic + SFTP \\
    E4$\times$C & 20 & 1.33 & 1.33 & 1.35 & MinIO (S3) \\
    \bottomrule
  \end{tabularx}
\end{table}

La restauración desde MinIO (estrategia~C) es sistemáticamente más rápida que desde
SFTP (estrategia~B): mediana de 1.34~s frente a 2.22~s, una reducción del
39.4\,\%. La separación entre ambas distribuciones es total ($\delta$~de Cliff~$=1.0$
en los tres escenarios), con significación extrema. La diferencia proviene del
protocolo de acceso al \emph{backend} —S3 frente a SFTP— y no del mecanismo de
inmutabilidad.

\section{Análisis estadístico — H1}
\label{sec:estadistica-h1}

La hipótesis~H1 se planteó originalmente sobre el \gls{mttr} extremo a extremo. Medido
así, el resultado es mixto y ruidoso (tabla~\ref{tab:mann-whitney-mttr}): en E1 la
estrategia~C es significativamente más rápida ($p = 0.001$), pero en E3 y E4 no hay
diferencia significativa, e incluso la mediana de C es mayor por efecto de la
latencia de detección incorporada al MTTR. El \gls{mttr} total, por tanto, no separa
limpiamente las estrategias.

\begin{table}[ht]
  \centering
  \caption{Mann-Whitney U y delta de Cliff para el \gls{mttr} total B~vs.~C.}
  \label{tab:mann-whitney-mttr}
  \rowcolors{2}{tabla-par}{white}
  \begin{tabularx}{\textwidth}{C{1.8cm} S[table-format=2.1] S[table-format=2.1] S[table-format=3.0] S[table-format=1.3] S[table-format=-1.2] X}
    \toprule
    \rowcolor{tabla-cabecera}
    \textcolor{white}{\textbf{Escenario}} &
    \textcolor{white}{\textbf{Med. B (s)}} &
    \textcolor{white}{\textbf{Med. C (s)}} &
    \textcolor{white}{\textbf{U}} &
    \textcolor{white}{\textbf{$p$-valor}} &
    \textcolor{white}{\textbf{Cliff $\delta$}} &
    \textcolor{white}{\textbf{Efecto}} \\
    \midrule
    E1 & 12.8 & 12.1 & 318 & 0.001 &  0.59 & grande \\
    E2 & {---}& 16.8 & {---} & {---} & {---} & B: impossible \\
    E3 &  7.9 & 12.1 & 180 & 0.598 & -0.10 & negligible \\
    E4 & 10.7 & 14.8 & 154 & 0.218 & -0.23 & pequeño \\
    \bottomrule
  \end{tabularx}
\end{table}

La comparación relevante es la de la \emph{velocidad de restauración pura}
(sección~\ref{sec:restauracion-pura}), que aísla el efecto del \emph{backend}. La
tabla~\ref{tab:mann-whitney} recoge el contraste B~vs.~C de esa métrica por escenario
y agregado.

\begin{table}[ht]
  \centering
  \caption{Mann-Whitney U, delta de Cliff e intervalo \textit{bootstrap} al 95\,\%
           para la velocidad de restauración B~vs.~C.}
  \label{tab:mann-whitney}
  \rowcolors{2}{tabla-par}{white}
  \begin{tabularx}{\textwidth}{C{1.6cm} S[table-format=1.2] S[table-format=1.2] S[table-format=4.0] c S[table-format=1.2] X}
    \toprule
    \rowcolor{tabla-cabecera}
    \textcolor{white}{\textbf{Escenario}} &
    \textcolor{white}{\textbf{Med. B}} &
    \textcolor{white}{\textbf{Med. C}} &
    \textcolor{white}{\textbf{U}} &
    \textcolor{white}{\textbf{$p$-valor}} &
    \textcolor{white}{\textbf{Cliff $\delta$}} &
    \textcolor{white}{\textbf{IC 95\,\% (B$-$C)}} \\
    \midrule
    E1              & 2.20 & 1.35 & 400 & $<10^{-7}$ & 1.00 & [0.78, 0.87] \\
    E3              & 2.22 & 1.34 & 400 & $<10^{-7}$ & 1.00 & [0.86, 0.90] \\
    E4              & 2.22 & 1.33 & 400 & $<10^{-7}$ & 1.00 & [0.87, 0.90] \\
    Agregado        & 2.22 & 1.34 & 3600 & $<10^{-20}$ & 1.00 & [0.86, 0.88] \\
    \bottomrule
  \end{tabularx}
\end{table}

\textbf{Reformulada sobre la velocidad de restauración, la hipótesis~H1 se confirma}:
la estrategia con almacenamiento inmutable restaura un 39.4\,\% más rápido que la
convencional, con un tamaño del efecto máximo ($\delta = 1.0$) e intervalos de
confianza que no contienen el cero en ningún escenario. Conviene subrayar que esta
ventaja proviene del protocolo S3 frente a SFTP, no de Object Lock en sí; la
aportación de Object Lock es de disponibilidad, no de velocidad, como muestra el
análisis de H3.

\section{Análisis estadístico — H2}
\label{sec:estadistica-h2}

\textbf{La hipótesis H2 se confirma}. Los tiempos de detección medianos oscilan entre
0.4 y 11.2~segundos, en todos los casos muy por debajo del umbral de 60~s. La
instrumentación con un \gls{siem} permite alertar sobre un ataque en curso en
cuestión de segundos, en contraste con los tiempos de permanencia (\emph{dwell time})
de meses que la industria reporta de forma recurrente para incidentes no
instrumentados. Conviene precisar que ambas cifras no son directamente comparables
—el \emph{dwell time} mide el intervalo desde la intrusión inicial hasta su
descubrimiento, mientras que $\tdetect$ mide el tiempo hasta alertar sobre una fase
de ataque ya iniciada y con reglas diseñadas para ella—, pero el contraste ilustra la
transformación cualitativa que aporta la detección automatizada.

Dentro del laboratorio, $\tdetect$ presenta baja variabilidad y depende del escenario
y de las reglas, no de la estrategia de \emph{backup}, lo que es coherente con el
diseño.

\section{Análisis estadístico — H3}
\label{sec:estadistica-h3}

\begin{table}[ht]
  \centering
  \caption{Test exacto de Fisher para la comparación de tasas de recuperación
           exitosa B~vs.~C.}
  \label{tab:fisher}
  \rowcolors{2}{tabla-par}{white}
  \begin{tabularx}{\textwidth}{X C{2cm} C{2cm} C{2.5cm} C{1.5cm}}
    \toprule
    \rowcolor{tabla-cabecera}
    \textcolor{white}{\textbf{Comparación}} &
    \textcolor{white}{\textbf{B (éxito)}} &
    \textcolor{white}{\textbf{C (éxito)}} &
    \textcolor{white}{\textbf{$p$-valor}} &
    \textcolor{white}{\textbf{Sig.}} \\
    \midrule
    E2 --- Destrucción de copias          & 0/20  & 20/20 & $< 0.001$ & Sí \\
    Global (E1+E2, 40 ejec.\ por estr.)   & 20/40 & 40/40 & $< 0.001$ & Sí \\
    \bottomrule
  \end{tabularx}
\end{table}

El test exacto de Fisher arroja $p < 0.001$ tanto para la comparación del
escenario~E2 específicamente (0/20 vs. 20/20) como para la comparación global sobre
los escenarios con pérdida de datos (20/40 vs. 40/40). \textbf{La hipótesis H3 se
confirma con evidencia estadística muy sólida}: la combinación de Object Lock
COMPLIANCE y mínimo privilegio \gls{iam} proporciona una ventaja cualitativa —no solo
cuantitativa— sobre el \emph{backup} convencional. La magnitud de la diferencia en E2
es la máxima posible: 0\,\% vs. 100\,\% de tasa de éxito.

\section{Síntesis narrativa por escenario}
\label{sec:sintesis-escenarios}

Esta sección reorganiza los resultados por escenario, con el objetivo de facilitar una
lectura orientada a la pregunta que un responsable de seguridad se plantearía ante
cada táctica de ataque.

En \textbf{E1 (cifrado local, T1486)}, las tres estrategias detectan el cifrado de
forma casi inmediata (mediana~$\tdetect \approx 0.4$~s, gracias a la regla de
\gls{fim} sobre la creación masiva de ficheros con extensión de cifrado). Para las
estrategias con \emph{backup}, el tiempo total de recuperación es similar
(mediana~MTTR 12.1--12.8~s), si bien la restauración pura de C es más rápida
(1.35~s frente a 2.20~s). Es el escenario donde el repositorio de \emph{backup} nunca
es objetivo directo, por lo que ambas estrategias recuperan el 100\,\% de los casos.

En \textbf{E2 (destrucción de copias, T1490)} se concentra la evidencia central del
trabajo. La estrategia~B no tiene ningún valor de MTTR calculable porque las veinte
ejecuciones resultan en \texttt{recovery\_status = impossible}: el atacante borra el
repositorio SFTP y la recuperación es inviable. La estrategia~C completa las veinte
recuperaciones (mediana~MTTR 16.8~s). La detección es más lenta que en E1
(mediana~$\tdetect \approx 11$~s en A y C) porque, al no existir repositorio SFTP que
borrar, la alerta procede del evento que confirma el bloqueo de Object Lock tras el
intento de borrado del \emph{bucket} protegido; en la estrategia~B, el borrado masivo
del repositorio SFTP se detecta en 0.7~s.

En \textbf{E3 (movimiento lateral, T1021.004)} la detección se produce en torno a
1.5~s mediante la regla de sesión SSH aceptada con origen en el \emph{fileserver}
(regla~100220), sin depender de umbrales de frecuencia. El \gls{mttr} total muestra
aquí a la estrategia~B con una mediana menor que la~C (7.9~s frente a 12.1~s), pero
esta diferencia es un artefacto de la cuantización del sondeo de detección y no es
significativa; la restauración pura, en cambio, es más rápida en~C.

En \textbf{E4 (exfiltración, T1041)} la detección es rápida (mediana~$\tdetect
\approx 1.2$~s), combinando la regla de sesión SSH con origen en el atacante y la de
creación del archivo comprimido en \texttt{/tmp}. Es un cambio sustancial respecto a
iteraciones previas del laboratorio, en las que este escenario presentaba los tiempos
de detección más altos por un artefacto en el canal de exfiltración.

\section{Verificación de integridad}

En las 140~ejecuciones con \texttt{recovery\_status = success} (60~de estrategia~B y
80~de estrategia~C), la variable \texttt{integrity\_match} toma el valor
\texttt{true} en todos los casos. No se observó ninguna restauración exitosa con
integridad fallida, lo que confirma que \restic garantiza la integridad byte a byte
de los datos restaurados en las condiciones del experimento.

\section{Resumen del contraste de hipótesis}

\needspace{12\baselineskip}
\begin{table}[ht]
  \centering
  \caption{Resumen del contraste de hipótesis del experimento.}
  \label{tab:hipotesis}
  \rowcolors{2}{tabla-par}{white}
  \begin{tabularx}{\textwidth}{C{1.5cm} X C{2.8cm} C{2.5cm}}
    \toprule
    \rowcolor{tabla-cabecera}
    \textcolor{white}{\textbf{Hip.}} &
    \textcolor{white}{\textbf{Enunciado}} &
    \textcolor{white}{\textbf{Resultado}} &
    \textcolor{white}{\textbf{Evidencia}} \\
    \midrule
    H1 & C reduce el tiempo de restauración vs.\ B &
        Confirmada (velocidad de restauración) & $\delta=1.0$, $p<10^{-20}$ \\
    H2 & $\tdetect < 60$~s mediante SIEM &
        Confirmada & Mediana 0.4--11.2~s \\
    H3 & C $>$ tasa éxito B ante destrucción &
        Confirmada & Fisher $p < 0.001$ \\
    \bottomrule
  \end{tabularx}
\end{table}

\end{document}
